\documentclass[11pt]{article}

\usepackage[margin=1in]{geometry}
\usepackage{amsmath,amssymb}
\usepackage{booktabs}
\usepackage{enumitem}
\usepackage[most]{tcolorbox}
\usepackage{xcolor}
\usepackage[hidelinks]{hyperref}

\definecolor{barcsblue}{HTML}{0072B2}
\newtcolorbox{keyidea}{colback=barcsblue!6,colframe=barcsblue,
  fonttitle=\bfseries,title=Key idea,boxrule=0.6pt,arc=2pt}
\newtcolorbox{example}[1][]{colback=gray!5,colframe=gray!60,
  fonttitle=\bfseries,title=Worked example#1,boxrule=0.6pt,arc=2pt}
\newtcolorbox{incode}{colback=white,colframe=gray!40,
  fonttitle=\bfseries,title=In the package,boxrule=0.4pt,arc=2pt}

\newcommand{\logit}{\operatorname{logit}}
\newcommand{\Var}{\operatorname{Var}}
\newcommand{\E}{\operatorname{E}}

\newlist{choices}{enumerate}{1}
\setlist[choices]{label=(\Alph*),leftmargin=2.2em,itemsep=1pt,topsep=2pt}

\title{\textbf{Understanding BARCS}\\[4pt]
\large A self-study guide to the methodology, with worked examples,\\
a multiple-choice quiz, and an answer key}
\author{Prepared for Hyun-Hwan Jeong and Kyu-Won Lee}
\date{BARCS version 0.2.0}

\begin{document}
\sloppy
\maketitle
\tableofcontents
\newpage

%=====================================================================
\section{The problem BARCS solves}
%=====================================================================

A pooled CRISPR screen produces, for every guide $g$ and every sequenced
library $i$, a read count $K_{gi}$.  Each library also has a total number of
mapped guide reads, $N_i$.  The question of interest is how guide abundance
depends on the experimental design: time, an ordered FACS bin, a dose, a
donor, a treatment, or an interaction between them.

CB\textsuperscript{2} answered this for two groups.  BARCS keeps the same
sampling model but replaces the two group means with a regression on an
arbitrary design matrix.  Three objects travel together:

\begin{center}
\begin{tabular}{lll}
\toprule
Object & Shape & Rule\\
\midrule
\texttt{counts} & guides $\times$ libraries & raw integer read counts\\
\texttt{totals} & one per library & unfiltered mapped-guide totals $N_i$\\
\texttt{data} & one row per library & same order as the columns of \texttt{counts}\\
\bottomrule
\end{tabular}
\end{center}

\begin{keyidea}
BARCS treats a guide's count as a proportion of its library.  The denominator
$N_i$ is the total recorded \emph{before} any guide is filtered out, because
removing guides and recomputing totals silently changes what every remaining
coefficient means.
\end{keyidea}

%=====================================================================
\section{The sampling model}
%=====================================================================

\subsection{From binomial to beta-binomial}

If every library of a given condition contained exactly the same guide
proportion $\mu_i$, sequencing would be pure binomial sampling:
$K_{gi}\sim\mathrm{Binomial}(N_i,\mu_i)$.  Real libraries are biological
replicates, so the true proportion itself varies from library to library.
BARCS lets the per-library proportion $P_i$ follow a beta distribution:
\begin{align}
K_i \mid P_i &\sim \mathrm{Binomial}(N_i, P_i), &
P_i &\sim \mathrm{Beta}\big(\mu_i\kappa,\,(1-\mu_i)\kappa\big), &
\rho &= \frac{1}{\kappa+1}.
\end{align}
Marginally, $K_i$ is beta-binomial with mean $N_i\mu_i$ and variance
\begin{equation}
\Var(K_i) = N_i\mu_i(1-\mu_i)\,\{1+(N_i-1)\rho\}.
\label{eq:var}
\end{equation}
The intraclass correlation $\rho\ge 0$ is the overdispersion: $\rho=0$
recovers the binomial.

\subsection{Why this separates depth from replication}

Divide by $N_i$ to get the variance of the observed proportion
$Y_i=K_i/N_i$:
\begin{equation}
\Var(Y_i)=\mu_i(1-\mu_i)\left\{\frac{1-\rho}{N_i}+\rho\right\}.
\label{eq:propvar}
\end{equation}
The first term shrinks as sequencing depth $N_i$ grows; the second does not.
Once a library is sequenced deeply, the remaining uncertainty is biological.

\begin{example}[: deeper sequencing barely helps]
Take $\mu=10^{-4}$ and $\rho=10^{-5}$.  The bracket in
Equation~\eqref{eq:propvar} is
\[
N=10^6:\ 10^{-6}+10^{-5}=1.1\times10^{-5},\qquad
N=10^7:\ 10^{-7}+10^{-5}=1.01\times10^{-5}.
\]
Ten times more reads reduce the variance by only about 8\%.  Adding a second
\emph{library}, by contrast, halves the variance of a mean.  This is the
whole reason BARCS takes its degrees of freedom from libraries, not reads.
\end{example}

%=====================================================================
\section{The regression model}
%=====================================================================

For guide $g$ in library $i$,
\begin{equation}
K_{gi}\mid N_i\sim\mathrm{BetaBinomial}(N_i,\mu_{gi},\rho_g),\qquad
\logit(\mu_{gi})=\mathbf{x}_i^\top\boldsymbol\beta_g .
\label{eq:model}
\end{equation}
Every guide has its own coefficients $\boldsymbol\beta_g$ and its own
dispersion $\rho_g$.  The design vector $\mathbf{x}_i$ is an ordinary row of
an R model matrix, so \texttt{\textasciitilde{} dose + batch},
\texttt{\textasciitilde{} replicate + time}, or
\texttt{\textasciitilde{} knockout * treatment} all work.

\paragraph{Interpreting a coefficient.}
$\beta_{gd}$ is the change in log-odds of the guide's library share per unit
of predictor $d$, holding the other columns fixed.  Because guide proportions
are tiny, $\logit(\mu)\approx\log\mu$, so $\beta_{gd}$ is approximately a log
fold change in abundance per unit of $x_d$.

\paragraph{Contrasts.}
Any linear combination $\mathbf{c}^\top\boldsymbol\beta$ can be tested, for
example the difference between two treatment arms
(\texttt{bb\_contrast()}).

%=====================================================================
\section{Estimation: feasible IRLS}
%=====================================================================

BARCS fits each guide separately by feasible iteratively reweighted least
squares.  Suppressing the guide index, with current $\boldsymbol\beta$ and
$\rho$:
\begin{align}
\eta_i&=\mathbf{x}_i^\top\boldsymbol\beta, &
\mu_i&=\frac{e^{\eta_i}}{1+e^{\eta_i}},\\
z_i&=\eta_i+\frac{Y_i-\mu_i}{\mu_i(1-\mu_i)}, &
w_i&=\frac{N_i\mu_i(1-\mu_i)}{1+(N_i-1)\rho}.
\end{align}
The update is weighted least squares on the working response:
$\widehat{\boldsymbol\beta}_{\text{new}}=(X^\top WX)^{-1}X^\top W\mathbf{z}$.

\paragraph{Estimating $\rho$.}
At each fitted mean, $\rho$ solves the Pearson estimating equation
\begin{equation}
\sum_{i=1}^{m}\frac{(K_i-N_i\hat\mu_i)^2}
{N_i\hat\mu_i(1-\hat\mu_i)\{1+(N_i-1)\hat\rho\}}=m-q,
\label{eq:pearson}
\end{equation}
where $m$ is the number of libraries and $q$ the number of coefficients.  The
left side decreases in $\rho$.  Three cases arise:
\begin{enumerate}[itemsep=1pt]
\item If the binomial Pearson statistic ($\rho=0$) is already $\le m-q$, the
counts are not overdispersed and $\hat\rho$ is truncated at 0.
\item Otherwise BARCS finds the root numerically.
\item If even $\rho$ at its upper bound leaves the statistic above $m-q$,
$\hat\rho$ is set to the bound and the covariance is multiplied by the
remaining Pearson scale.
\end{enumerate}
Coefficient and dispersion updates alternate until the coefficients change by
less than the tolerance.

\paragraph{The weight saturates.}
For $\rho>0$, $w_i\to\mu_i(1-\mu_i)/\rho$ as $N_i\to\infty$.  A library's
weight cannot grow without bound with depth, which is Equation~\eqref{eq:propvar}
seen from the fitting side.

\begin{incode}
\texttt{bbreg()} fits one guide; \texttt{bb\_screen()} fits every row of the
count matrix (optionally in parallel).  The weighted cross-products are
computed in C++ with RcppArmadillo.
\end{incode}

%=====================================================================
\section{Inference for a single guide}
%=====================================================================

At convergence the covariance is $\widehat{\mathrm{Cov}}(\widehat{\boldsymbol\beta})
=(X^\top\widehat WX)^{-1}$ and, for coefficient $j$,
\begin{equation}
t_j=\frac{\hat\beta_j}{\sqrt{\widehat{\mathrm{Cov}}_{jj}}}\ \sim\ t_{m-q}.
\end{equation}
The reference degrees of freedom are $m-q$: independent libraries minus
coefficients.  Total read count never enters.

\begin{example}[: degrees of freedom in the Cas13 screens]
Each cell line has 6 libraries (days 0, 7, 14 $\times$ 2 replicates) and the
model \texttt{\textasciitilde{} replicate + time} has 3 coefficients
(intercept, replicate, time).  Each guide therefore has $6-3=3$ residual
degrees of freedom.  A variance estimated from 3 degrees of freedom is very
noisy, which motivates the next section.
\end{example}

The $t$ reference is heavier-tailed than a normal, but it treats
$\hat\rho$ as known.  It is not an exact correction for dispersion
uncertainty; moderation and negative controls are the practical safeguards.

%=====================================================================
\section{Dispersion moderation (on by default)}
%=====================================================================

\subsection{The idea}

With 3 residual degrees of freedom, some guides look quiet by chance and get
tiny standard errors, others look noisy and get huge ones.  Moderation shrinks
each guide's variance toward a library-wide trend, as limma and edgeR do.

\subsection{The procedure}

\begin{enumerate}[itemsep=2pt]
\item \textbf{Raw inflation.} Let $\nu=m-q$ and
$s_g^2=Q^{(0)}_g/\nu$, where $Q^{(0)}_g$ is the guide's Pearson statistic
under the \emph{binomial} model.  Unlike $\hat\rho$, $s_g^2$ is not
truncated, so it behaves like a scaled $\chi^2_\nu/\nu$.
\item \textbf{Trend.} Fit a lowess curve of $\log s_g^2$ against log mean
CPM (span 0.5).  Call the fitted value $\log s^2_{\text{trend},g}$.
\item \textbf{Prior degrees of freedom (Smyth 2004).} Let
$e_g=\log s_g^2-\log s^2_{\text{trend},g}-\psi(\nu/2)+\log(\nu/2)$, with
$\psi$ the digamma function.  Compute
\[
\text{excess}=\widehat{\Var}(e_g)-\psi'(\nu/2).
\]
If $\text{excess}>0$, solve $\psi'(d_0/2)=\text{excess}$ for $d_0$;
otherwise $d_0=\infty$ (complete shrinkage).
\item \textbf{Prior scale.} $s_{0g}^2=\exp\{\log s^2_{\text{trend},g}
+\bar e+\psi(d_0/2)-\log(d_0/2)\}$.
\item \textbf{Posterior.}
\begin{equation}
\tilde s_g^2=\frac{\nu s_g^2+d_0 s_{0g}^2}{\nu+d_0}.
\end{equation}
\item \textbf{Rescale.} The fitted standard error already reflects
$\max(1,s_g^2)$, so
\[
\mathrm{SE}_{g,\text{mod}}=\mathrm{SE}_g\sqrt{\tilde s_g^2/\max(1,s_g^2)},\qquad
t_{g,\text{mod}}=\hat\beta_g/\mathrm{SE}_{g,\text{mod}},\qquad
\text{df}=\nu+d_0.
\]
\end{enumerate}
The coefficient $\hat\beta_g$ never changes; only its standard error and
reference distribution do.

\begin{example}[: one guide]
$\nu=3$, $s_g^2=4$, prior $s_{0g}^2=1.5$, $d_0=12$.
\[
\tilde s_g^2=\frac{3(4)+12(1.5)}{15}=2.0,\qquad
\text{rescale}=\sqrt{2/4}=0.71,\qquad \text{df}=15.
\]
The noisy guide's standard error shrinks by 29\% and its test now uses a
$t_{15}$ reference instead of $t_3$.  A guide with $s_g^2=0.5$ would instead
be pulled \emph{up} toward the prior, protecting against false confidence.
\end{example}

\begin{incode}
\texttt{bb\_screen(moderate = NULL)} moderates when at least 50 guides are
usable; \texttt{moderate = FALSE} turns it off.  Options:
\texttt{trend = FALSE} shrinks to one constant; \texttt{one\_way = TRUE} only
ever raises a variance; \texttt{borrow\_df = FALSE} keeps $\nu$ degrees of
freedom.  A result cannot be moderated twice.
\end{incode}

\paragraph{Why it matters.} In the Cas13 screens, moderation raised
essential-gene average precision from 0.838 to 0.876, level with limma--voom
and edgeR-QL, and improved every cell line.

%=====================================================================
\section{The denominator: what a coefficient is relative to}
%=====================================================================

\subsection{Full-library totals and composition shift}

With $N_i$ as the denominator, $\beta$ measures change in library share.  If
many guides deplete (a dropout screen), every surviving guide's share rises
even when its absolute abundance is unchanged.  Null guides then have
positive coefficients, and their $t$ statistics can be large.

\subsection{Control-based totals}

\texttt{barcs\_control\_totals()} holds a chosen control class $C$ at a
constant share of every library:
\begin{equation}
N_i^{\text{ctrl}}=\operatorname{round}\!\left(
\frac{\sum_{g\in C}K_{gi}}{\overline{\sum_{g\in C}K_{g\cdot}}}\;
\overline{N_\cdot}\right),
\quad\text{raised if needed so that } N_i^{\text{ctrl}}\ge\max_g K_{gi}.
\end{equation}
Coefficients then describe abundance relative to the controls.

\begin{example}[: control totals]
Two libraries; two control guides with 500 reads each in both; one hit with
1,000 reads in library 1 and 100 in library 2.  Full totals are 2,000 and
1,100 (mean 1,550).  Control totals are $1000\times1550/1000=1550$ in both
libraries, so the controls keep the same share and the hit's depletion is
measured against them.
\end{example}

\paragraph{Choosing the control class.} Non-targeting guides do not cut;
safe-harbor guides cut but hit no gene.  Normalizing on safe-harbor guides
absorbs the shared cutting response.  Guides used to build the totals should
not also be used to judge the result.

%=====================================================================
\section{Negative-control calibration}
%=====================================================================

When the model's assumptions are violated (for example, FACS bins from one
donor partition the same cells), null $t$ statistics are too wide.
\texttt{bb\_calibrate\_controls()} rescales all $t$ statistics by one factor
estimated from negative controls:
\begin{itemize}[itemsep=1pt]
\item \texttt{tail\_quantile}: $\text{scale}=Q_{0.95}(|t_{\text{ctrl}}|)/t_{\nu,0.975}$;
\item \texttt{qq\_slope}: the slope, through the origin, of control quantiles
against $t$ quantiles over the 0.50--0.95 band (less variable);
\end{itemize}
and then $\text{scale}\leftarrow\max(1,\text{scale})$, so calibration can only
make tests more conservative.  Standard errors are multiplied and $t$
statistics divided by the scale; estimates are unchanged.

\begin{example}[: tail scale]
Controls' 95th percentile of $|t|$ is 3.5 with $\nu=4$.  Since
$t_{4,0.975}=2.776$, the scale is $3.5/2.776=1.26$: every $t$ is divided by
1.26.
\end{example}

\paragraph{Honest evaluation.} If the same controls choose the scale and are
then used to report the null rate, the 5\% rate is guaranteed by construction.
The IL2RA analysis therefore uses five-fold cross-fitting: each control is
evaluated with a scale estimated from the other four folds.

%=====================================================================
\section{From guides to genes}
%=====================================================================

\subsection{Directional Stouffer}

For the $j$th valid guide of gene $g$, convert the two-sided $p$-value back to
a signed normal score, then sum:
\begin{equation}
z_{gj}=\operatorname{sign}(\hat\beta_{gj})\,\Phi^{-1}\!\left(1-\tfrac{p_{gj}}{2}\right),
\qquad
Z_g=\frac{\sum_{j=1}^{m_g}z_{gj}}{\sqrt{m_g}},\qquad p_g=2\Phi(-|Z_g|).
\end{equation}
Concordant guides reinforce each other; discordant ones cancel.  The reported
gene effect is the \emph{median} guide coefficient, which resists one extreme
guide.  Gene $p$-values are then BH-adjusted.

\subsection{The independence problem}

$\sqrt{m_g}$ is the standard deviation of a sum of $m_g$ \emph{independent}
standard normals.  But guides of one gene share every library; if they share
noise, their null scores have positive correlation $r$, and
\begin{equation}
\Var\Big(\sum_j z_{gj}\Big)=m_g+m_g(m_g-1)\,r .
\end{equation}
BARCS divides by the square root of this instead:
\begin{equation}
Z_g^{\text{c}}=\frac{\sum_{j}z_{gj}}{\sqrt{m_g+m_g(m_g-1)\hat r}} .
\end{equation}

\begin{example}[: what correlation does]
Five guides, each with $z=2.5$, so $\sum z=12.5$.  Uncorrected,
$Z=12.5/\sqrt5=5.59$.  Suppose instead the five guides have sum
$\sum z=5.59$ (each $z\approx1.12$).  Uncorrected $Z=2.5$, $p=0.012$.
With $r=0.4$, the variance is $5+5\cdot4\cdot0.4=13$ and
$Z^{\text{c}}=5.59/\sqrt{13}=1.55$, $p=0.12$: no longer significant.
\end{example}

\subsection{Estimating $r$ without being fooled by real signal}

\begin{enumerate}[itemsep=2pt]
\item Take each guide's Pearson residual vector across libraries and scale it
to unit length.
\item For each gene with at least two guides, average the pairwise inner
products (correlations); then average over genes, giving each gene equal
weight (so a large control ``gene'' cannot dominate).
\item Do the same for deterministic pairs of guides from \emph{different}
genes; this baseline removes correlation shared by the whole library.
\item $\hat r=\min\{1,\max(0,\ \text{within}-\text{between})\}$.
\end{enumerate}

\begin{keyidea}
Residuals are what is left after the fitted mean has absorbed every design
effect, including a true gene effect.  A gene whose guides all deplete
strongly therefore does \emph{not} have correlated residuals; only noise
shared by its guides does.  That is why the correction does not penalize real
hits.
\end{keyidea}

In simulations with $r=0.4$, the corrected gene-level type I error was
0.018--0.052 versus up to 0.262 uncorrected.  In the real screens analyzed,
$\hat r$ was 0.008--0.026, so the correction changed little there.

\subsection{Experimental alternatives}

Four other summaries pool effect sizes rather than tests:
\texttt{bb\_gene\_normal()}, \texttt{bb\_gene\_consistency()},
\texttt{bb\_gene\_partial\_pool()}, and
\texttt{bb\_gene\_eb\_moderate()}.  They ranked genes worse than Stouffer in
the simulated FACS benchmarks and are kept as sensitivity analyses.

%=====================================================================
\section{How the benchmarks measure success}
%=====================================================================

\begin{center}
\begin{tabular}{lp{0.62\textwidth}}
\toprule
Metric & Meaning\\
\midrule
AUROC & Probability a random positive outranks a random negative.\\
Average precision (AP) & Mean precision at each true positive's rank;
sensitive to the top of the list.\\
Recall at 5\% null FPR & Fraction of positives above the score that only
5\% of null genes exceed; threshold-free with respect to calibration.\\
Recall at FDR 0.10 & Fraction of positives called at the method's own
nominal FDR; depends on calibration.\\
Realized FDP & Fraction of calls that are false, when truth is known.\\
Proxy null & Genes assumed null (e.g. unexpressed lncRNAs); imperfect, so
null rates are diagnostics.\\
\bottomrule
\end{tabular}
\end{center}

\paragraph{Headline results (BARCS 0.2.0).}
\begin{center}
\begin{tabular}{lll}
\toprule
Benchmark & BARCS & Comparison\\
\midrule
Cas13 AP & 0.876 & limma 0.876, edgeR 0.874, MAGeCK 0.877\\
Cas13 recall @ FDR 0.10 & 0.767 & 0.692--0.767\\
IL2RA validated regulators & 23/26, 61 calls & MAGeCK 17/26, 72 calls\\
CRISPulator F1 (moderated) & 0.847 & unmoderated 0.811\\
Null grid, correlated guides & 0.018--0.052 & uncorrected up to 0.262\\
A375 recall @ FDR 0.05 & 0.883 & MAGeCK 0.710\\
\bottomrule
\end{tabular}
\end{center}

%=====================================================================
\section{Limitations to keep in mind}
%=====================================================================

\begin{itemize}[itemsep=2pt]
\item \textbf{Independent libraries.} FACS bins from one donor and repeated
harvests from one culture are dependent.  Donor indicators adjust means, and
controls fix the miscalibration, but the dependence itself is not modeled.
\item \textbf{Plug-in dispersion.} The $t$ reference treats $\hat\rho$ as
known; moderation reduces, but does not remove, this approximation.
\item \textbf{Likelihood-compatible counts.} The Cas13 benchmark used
normalized, batch-corrected values rounded to integers; a raw-count rerun
would be the stronger test.
\item \textbf{Denominator choice.} Full-library and control-based
coefficients answer different questions; choose before looking at results.
\end{itemize}

%=====================================================================
\section{Pipeline at a glance}
%=====================================================================

\begin{enumerate}[itemsep=2pt]
\item \texttt{barcs\_quant()}: FASTQ $\rightarrow$ \texttt{counts},
\texttt{totals}.
\item Optional \texttt{barcs\_control\_totals()}: choose the reference.
\item \texttt{bb\_screen()}: per-guide beta-binomial regression, $t$ tests,
dispersion moderation, residual correlation $\hat r$.
\item Optional \texttt{bb\_calibrate\_controls()}: negative-control scale.
\item \texttt{bb\_gene\_stouffer()}: correlation-aware gene statistic,
median effect, BH FDR.
\end{enumerate}

%=====================================================================
\newpage
\section{Quiz}
%=====================================================================

Choose the single best answer.  The answer key with explanations follows on a
separate page.

\begin{enumerate}[label=\textbf{Q\arabic*.},leftmargin=2.6em,itemsep=8pt]

\item In BARCS, the total $N_i$ for library $i$ should be
\begin{choices}
\item the median count across guides
\item the sum of counts after removing low-count guides
\item the full mapped-guide total recorded before any guide filtering
\item the number of FASTQ reads in the file
\end{choices}

\item The beta-binomial variance of a guide count is
$N\mu(1-\mu)\{1+(N-1)\rho\}$.  Setting $\rho=0$ gives
\begin{choices}
\item zero variance
\item a binomial variance
\item a Poisson variance
\item a negative-binomial variance
\end{choices}

\item With $\mu=10^{-4}$ and $\rho=10^{-5}$, increasing depth from $10^6$ to
$10^7$ reads reduces the variance of the observed proportion by roughly
\begin{choices}
\item 8\%
\item 50\%
\item 90\%
\item 0\%
\end{choices}

\item The residual degrees of freedom of a BARCS guide test are
\begin{choices}
\item number of independent libraries minus number of coefficients
\item number of guides per gene minus one
\item fixed at 30
\item total reads minus number of coefficients
\end{choices}

\item In the model $\logit(\mu_{gi})=\mathbf x_i^\top\boldsymbol\beta_g$, a
coefficient $\beta_{gd}$ is best interpreted as
\begin{choices}
\item the overdispersion of the guide
\item the probability that the guide is essential
\item the change in raw read count per unit of $x_d$
\item approximately the log fold change in guide abundance per unit of
$x_d$, holding other columns fixed
\end{choices}

\item As $N_i\to\infty$ with $\rho>0$, the IRLS weight
$w_i=N_i\mu_i(1-\mu_i)/\{1+(N_i-1)\rho\}$
\begin{choices}
\item grows without bound
\item approaches zero
\item approaches $\mu_i(1-\mu_i)/\rho$
\item equals $N_i$
\end{choices}

\item If a guide's binomial Pearson statistic is already below $m-q$, BARCS
sets $\hat\rho$ to
\begin{choices}
\item 0
\item the library-wide median
\item $1/(m-q)$
\item its upper bound
\end{choices}

\item A Cas13 cell line has 6 libraries and the model
\texttt{\textasciitilde{} replicate + time}.  How many residual degrees of
freedom does each guide have?
\begin{choices}
\item 5
\item 3
\item 4
\item 2
\end{choices}

\item Dispersion moderation changes
\begin{choices}
\item the library totals
\item the standard error and the reference degrees of freedom, but not
$\hat\beta_g$
\item the guide-to-gene assignment
\item the coefficient estimate $\hat\beta_g$
\end{choices}

\item Moderation shrinks each guide's log variance inflation toward
\begin{choices}
\item a lowess trend in log mean CPM across all guides
\item zero
\item the inflation of the negative controls only
\item the largest inflation in the library
\end{choices}

\item The prior degrees of freedom $d_0$ are estimated by
\begin{choices}
\item a fixed constant of 4
\item maximum likelihood on the counts
\item cross-validation on essential genes
\item the scaled-$F$ moment method of Smyth (2004), using digamma and
trigamma functions
\end{choices}

\item With $\nu=3$, $s_g^2=4$, $s_{0g}^2=1.5$, $d_0=12$, the moderated
inflation $\tilde s_g^2$ is
\begin{choices}
\item 2.75
\item 4.0
\item 2.0
\item 1.5
\end{choices}

\item Why is moderation based on $s_g^2=Q_g^{(0)}/\nu$ rather than
on $\hat\rho_g$?
\begin{choices}
\item $\hat\rho_g$ depends on read depth only
\item $s_g^2$ ignores the design
\item $\hat\rho_g$ is not stored by \texttt{bb\_screen()}
\item $\hat\rho_g$ is truncated at zero, so it does not behave like a
scaled $\chi^2$ quantity
\end{choices}

\item By default, \texttt{bb\_screen()} moderates dispersions when
\begin{choices}
\item the user supplies negative controls
\item the design has more than two groups
\item at least 50 guides are usable
\item always, even with 5 guides
\end{choices}

\item In a dropout screen with full-library totals, null guides tend to show
\begin{choices}
\item negative coefficients, because all guides deplete
\item undefined coefficients
\item coefficients centered at zero
\item positive coefficients, because surviving guides gain library share
\end{choices}

\item \texttt{barcs\_control\_totals()} produces totals that
\begin{choices}
\item are normalized to counts per million
\item hold a chosen control class at a constant share of every library
\item exclude the control guides
\item equal the number of FASTQ reads
\end{choices}

\item Normalizing on safe-harbor rather than non-targeting guides
\begin{choices}
\item absorbs the shared cutting response into the reference
\item removes all composition shift and all cutting effects
\item makes the totals non-integer
\item has no effect
\end{choices}

\item \texttt{bb\_calibrate\_controls()} with \texttt{min\_scale = 1} can
\begin{choices}
\item only make tests more liberal
\item only make tests more conservative (or leave them unchanged)
\item change the gene assignment
\item change coefficient estimates
\end{choices}

\item Control 95th percentile $|t|=3.5$ with $\nu=4$ ($t_{4,0.975}=2.776$).
The tail scale is about
\begin{choices}
\item 3.50
\item 1.00
\item 1.26
\item 0.79
\end{choices}

\item Why did the IL2RA analysis use five-fold cross-fitting of control
calibration?
\begin{choices}
\item Because the screen has five donors
\item To speed up computation
\item Because estimating the scale and evaluating the null rate on the same
controls makes a 5\% rate true by construction
\item To estimate guide correlation
\end{choices}

\item In directional Stouffer, a guide's signed score is
\begin{choices}
\item $-\log_{10}p_{gj}$
\item $\hat\beta_{gj}/\mathrm{SE}_{gj}$ always
\item $\operatorname{sign}(\hat\beta_{gj})\,\Phi^{-1}(1-p_{gj}/2)$
\item $\Phi^{-1}(p_{gj})$
\end{choices}

\item The reported BARCS gene effect is
\begin{choices}
\item the mean guide coefficient
\item the median guide coefficient
\item the Stouffer $Z$
\item the largest guide coefficient
\end{choices}

\item If $m$ null guide scores each have variance 1 and pairwise correlation
$r$, the variance of their sum is
\begin{choices}
\item $m$
\item $m+m(m-1)r$
\item $m^2r$
\item $m(1-r)$
\end{choices}

\item With 5 guides and $r=0.4$, the corrected Stouffer denominator is
$\sqrt{13}$ instead of $\sqrt5$.  A raw $Z=2.5$ ($p=0.012$) becomes about
\begin{choices}
\item $Z=4.0$, $p<0.001$
\item $Z=1.55$, $p=0.12$
\item $Z=2.5$, $p=0.012$
\item $Z=0$, $p=1$
\end{choices}

\item BARCS estimates within-gene correlation from residuals because
\begin{choices}
\item a true gene effect is absorbed by the fitted mean, so it does not
inflate residual correlation, whereas shared noise does
\item residuals do not depend on the design
\item residuals are always independent
\item residuals are smaller than counts
\end{choices}

\item Why does BARCS subtract the average correlation of guides from
\emph{different} genes?
\begin{choices}
\item To account for guide length
\item To make $\hat r$ negative
\item Because different genes are perfectly correlated
\item To remove correlation shared by the whole library, such as
normalization artifacts
\end{choices}

\item In the within-gene average, each gene receives equal weight so that
\begin{choices}
\item the estimate is always zero
\item genes with many guides dominate
\item essential genes are ignored
\item a large group of control guides labelled as one ``gene'' cannot
dominate the estimate
\end{choices}

\item In the real screens analyzed, the estimated within-gene correlation was
\begin{choices}
\item 0.008--0.026, so the correction changed little
\item larger than 0.5 in every screen
\item about 0.4
\item exactly zero by definition
\end{choices}

\item Recall at a 5\% empirical null false-positive rate differs from recall
at FDR 0.10 because it
\begin{choices}
\item ignores ranking
\item requires negative controls
\item uses only essential genes
\item does not depend on how well the method's $p$-values are calibrated
\end{choices}

\item In the Cas13 benchmark, the main reason BARCS caught up with limma and
edgeR in version 0.2.0 was
\begin{choices}
\item removing the day-7 libraries
\item switching to a negative-binomial likelihood
\item using control-based totals
\item dispersion moderation borrowing strength across guides
\end{choices}

\item Which statement about FACS bins from one donor is correct?
\begin{choices}
\item They should be summed before analysis
\item They require a binomial model
\item They partition the same cells and are dependent; donor indicators
adjust means but do not model that dependence
\item They are independent libraries, so BARCS models them exactly
\end{choices}

\item The experimental gene summaries (\texttt{bb\_gene\_partial\_pool()},
etc.) are kept because
\begin{choices}
\item they are required by \texttt{bb\_screen()}
\item they replace calibration
\item they ranked genes best in the benchmarks
\item they serve as sensitivity analyses; Stouffer ranked genes better in the
simulated FACS benchmarks
\end{choices}

\end{enumerate}

%=====================================================================
\newpage
\section{Answer key with explanations}
%=====================================================================

\begin{enumerate}[label=\textbf{Q\arabic*.},leftmargin=2.6em,itemsep=5pt]
\item \textbf{C.} The denominator must describe the library that was
sequenced. Recomputing it after filtering silently changes every remaining
guide's share (Section~1).
\item \textbf{B.} With $\rho=0$ the bracket is 1, leaving $N\mu(1-\mu)$, the
binomial variance.
\item \textbf{A.} The bracket $(1-\rho)/N+\rho$ goes from $1.1\times10^{-5}$
to $1.01\times10^{-5}$, about an 8\% reduction: biological variation
dominates.
\item \textbf{A.} $m-q$. Depth affects precision within a library but adds no
degrees of freedom.
\item \textbf{D.} Guide proportions are tiny, so log-odds $\approx$ log
proportion; the coefficient is a conditional log fold change.
\item \textbf{C.} The denominator grows like $N\rho$, so the weight saturates
at $\mu(1-\mu)/\rho$.
\item \textbf{A.} No overdispersion beyond binomial is detectable, so
$\hat\rho$ is truncated at 0.
\item \textbf{B.} $6-3=3$.
\item \textbf{B.} Moderation rescales the standard error and adds prior
degrees of freedom; $\hat\beta_g$ is unchanged.
\item \textbf{A.} The prior follows a lowess trend of log inflation against
log mean CPM (or a constant with \texttt{trend = FALSE}).
\item \textbf{D.} The excess variance of the centered log inflations over
$\psi'(\nu/2)$ is matched to $\psi'(d_0/2)$.
\item \textbf{C.} $(3\cdot4+12\cdot1.5)/15=30/15=2.0$.
\item \textbf{D.} Truncation at zero distorts the distribution; the
untruncated Pearson inflation behaves like a scaled $\chi^2_\nu/\nu$, which
the scaled-$F$ model assumes.
\item \textbf{C.} With fewer than 50 usable guides the prior cannot be
estimated reliably, so moderation is skipped unless forced.
\item \textbf{D.} Composition shift: as essential guides drop out, survivors
gain share, so null coefficients shift positive.
\item \textbf{B.} Control totals keep the control class at a fixed share,
rescaled to the mean library size and rounded.
\item \textbf{A.} Safe-harbor guides cut DNA, so normalizing on them removes
the shared cutting response from targeted effects.
\item \textbf{B.} The scale is floored at 1, so test statistics can only be
divided by a number $\ge1$.
\item \textbf{C.} $3.5/2.776=1.26$.
\item \textbf{C.} A held-out evaluation avoids the circularity of judging
calibration on the controls that set it.
\item \textbf{C.} The two-sided $p$ is mapped back to a normal deviate and
given the coefficient's sign.
\item \textbf{B.} The median resists a single extreme guide.
\item \textbf{B.} $m$ variance terms plus $m(m-1)$ covariance terms of size
$r$.
\item \textbf{B.} $Z^{\text c}=2.5\sqrt5/\sqrt{13}=1.55$, $p=0.12$.
\item \textbf{A.} Residuals exclude fitted design effects, so only shared
noise produces residual correlation.
\item \textbf{D.} The between-gene baseline removes library-wide
correlation that is not specific to a gene.
\item \textbf{D.} Equal gene weighting prevents a large pseudo-gene from
dominating the average.
\item \textbf{A.} Estimated $\hat r$ was 0.008--0.026 (Cas13), 0.011
(IL2RA), and 0.011 (Evers).
\item \textbf{D.} It fixes the false-positive rate empirically using proxy
nulls, so it compares ranking independent of $p$-value calibration.
\item \textbf{D.} Moderation raised AP from 0.838 to 0.876; without it each
guide's variance rests on only 3 degrees of freedom.
\item \textbf{C.} Bins partition one cell pool and are negatively dependent;
negative-control calibration corrects the resulting miscalibration, and
Waterbear models the dependence directly.
\item \textbf{D.} They are retained for sensitivity analysis; Stouffer gave
higher average precision and F1 in most simulated scenarios.
\end{enumerate}

\end{document}
